\documentclass{article}

\PassOptionsToPackage{numbers, sort, compress}{natbib}
\usepackage{neurips_2025}
\nolinenumbers

\usepackage[T1]{fontenc}
\usepackage[utf8]{inputenc}
\usepackage{hyperref}
\usepackage{url}
\usepackage{booktabs}
\usepackage{nicefrac}
\usepackage{microtype}
\usepackage{xcolor}
\usepackage{caption}

%%%

\usepackage{subcaption}
\usepackage{graphicx}
\usepackage{multirow}
\usepackage{amsmath,amssymb,amsfonts}
\usepackage{amsthm}
\usepackage{mathrsfs}
\usepackage{xcolor}
\usepackage{textcomp}
\usepackage{manyfoot}
\usepackage{algorithm}
\usepackage{algorithmicx}
\usepackage{algpseudocode}
\usepackage{listings}

\newtheorem{theorem}{Theorem} % continuous numbers
%%\newtheorem{theorem}{Theorem}[section] % sectionwise numbers
%% optional argument [theorem] produces theorem numbering sequence instead of independent numbers for Proposition
\newtheorem{proposition}[theorem]{Proposition}% 
\newtheorem{lemma}{Lemma}% 
%%\newtheorem{proposition}{Proposition} % to get separate numbers for theorem and proposition etc.

\newtheorem{example}{Example}
\newtheorem{remark}{Remark}

\newtheorem{definition}{Definition}
\newtheorem{assumption}{Assumption}

%%%

\newcommand{\R}{\mathbb{R}}
\newcommand{\cN}{\mathcal{N}}
\newcommand{\cL}{\mathcal{L}}
\newcommand{\cK}{\mathcal{K}}
\newcommand{\cD}{\mathcal{D}}
\newcommand{\cF}{\mathcal{F}}
\newcommand{\KL}{\mathrm{KL}}
\newcommand{\norm}[1]{\left\|#1\right\|}
\newcommand{\normsq}[1]{\left\|#1\right\|^{2}}
\newcommand{\ip}[2]{\langle #1,\, #2\rangle}
\newcommand{\hats}{\hat{s}}
\newcommand{\T}{\mathcal{T}}
\newcommand{\cG}{\mathcal{G}}


\title{TEMPO: Regime-Aware Operator Learning with Locally Adapted POD Bases}


% The \author macro works with any number of authors. There are two commands
% used to separate the names and addresses of multiple authors: \And and \AND.
%
% Using \And between authors leaves it to LaTeX to determine where to break the
% lines. Using \AND forces a line break at that point. So, if LaTeX puts 3 of 4
% authors names on the first line, and the last on the second line, try using
% \AND instead of \And before the third author name.





\begin{document}


\maketitle

\begin{abstract}

Operator learning methods typically approximate the solution map of a physical system with a single global model — an assumption that fails when solutions change character fundamentally across the parameter space.

We propose TEMPO: a framework that identifies latent dynamical regimes in snapshot data via Expectation-Maximization, builds a dedicated Neural-POD basis for each, and learns to route new inputs to the right regime at inference time. On multi-regime benchmarks, TEMPO achieves substantially lower errors than POD-DeepONet and vanilla DeepONet, with interpretable structure that reflects the underlying physics.

\end{abstract}

\textbf{Keywords:} Operator learning, DeepONet, Proper Orthogonal Decomposition, Expectation-Maximization, Gaussian mixture models, Regime discovery

\section{Introduction}\label{sec:intro}

Operator learning addresses the approximation of mappings between infinite-dimensional function spaces, arising naturally in the solution of parametric partial differential equations (PDEs) governing physical systems~\cite{Kovachki2024}. This paradigm has demonstrated substantial promise across both  theoretical analysis~\cite{Kovachki2024, Lu2019} and practical applications including weather modeling, quantum systems, fluid and sound dynamics~\cite{Pathak2022, Olivieri2024,  Dave2025, Tucker2025}.

Among the most prominent architectures, Deep Operator Networks (DeepONet)~\cite{Lu2019} approximate operators via paired branch and trunk networks, where the branch encodes the input function and the trunk encodes the query location.The Fourier Neural Operator (FNO)~\cite{Li2020} independently parameterizes the solution operator in the spectral domain, offering resolution-invariance and strong empirical performance. POD-DeepONet~\cite{Lu2022} improves upon standard DeepONet by replacing the learned trunk network with a fixed basis obtained from Proper Orthogonal Decomposition (POD) of training snapshots, yielding more compact and interpretable representations.

However, a single global POD basis is inherently limited for problems whose solutions exhibit qualitatively distinct dynamical regimes — such as solutions transitioning between laminar and turbulent behavior, or dispersive waves with parameter-dependent structure. In such settings, the optimal basis varies across the solution manifold, and a single linear subspace cannot represent the full solution diversity.

Several approaches enrich the POD basis to overcome this limitation. Sharma and Shankar~\cite{Sharma2025} update the DeepONet trunk with multiple complementary bases — precomputed POD modes and a partition-of-unity mixture-of-experts (PoU-MoE) trunk, trained jointly with a single branch. On 2D and 3D problems with sharp spatial gradients, the POD-PoU ensemble achieves up to 2$\times$ lower errors over POD-DeepONet~\cite{Lu2022} and 2--4$\times$ over vanilla DeepONet.


Moreover, the linear structure of the POD basis is insufficient for strongly nonlinear solution manifolds. Mou et al.~\cite{Mou2026} propose Neural-POD, which constructs nonlinear orthogonal basis functions via iterative residual minimization analogous to Gram--Schmidt orthogonalization, representing each  spatial mode and temporal component as neural networks trained sequentially on the residual variance.


Yet both approaches enrich the basis without identifying the underlying dynamical regimes present in the data. This idea has been pursued in the reduced order modeling (ROM) literature, where snapshot ensembles are partitioned into clusters and a separate local POD basis is constructed for each~\cite{Amsallem2012, Hess2019}. More recently, Daniel et al.~\cite{Daniel2020} extended this paradigm by clustering solution subspaces on the Grassmann manifold and training a deep neural network to select the appropriate local model from a dictionary at inference time. However, these methods operate offline on fixed snapshot sets, employ hard cluster assignments, and do not integrate regime discovery with operator learning in an end-to-end trainable architecture.

In this work, we propose \textbf{TEMPO} (\textit{Regime-Aware Operator Learning with Locally Adapted POD Bases}). Our method integrates regime discovery directly into operator learning: a Gaussian mixture model over Neural-POD coefficient space identifies latent dynamical regimes via soft EM, constructing a dedicated locally adapted basis for each. At inference time, a learned gating network selects among regimes, producing predictions that are accurate across the full parameter space.


\section{Problem Statement}

The behavior of many physical systems is governed by a set of parameters $\kappa \in \cK \subset \R^{d_\kappa}$: viscosity controls whether flow is laminar or turbulent, boundary forcing shapes wave patterns, material coefficients determine whether stress propagates smoothly or develops shocks. As $\kappa$ varies across the parameter space $\cK$, the corresponding solutions can change not just quantitatively but qualitatively, demanding fundamentally different representations in each regime. Predicting these solutions accurately and instantly, without re-running the underlying simulator, is the central problem we address.

Let $\Omega \subset \R^d$ be a bounded spatial domain and $\T \subset \R_{\geq 0}$ a time interval.

We are given a dataset of $N$ solution trajectories:
\begin{equation}
  \cD = \bigl\{\,(u^{(i)}_0,\,\kappa^{(i)},\,s^{(i)})\,\bigr\}_{i=1}^{N},
  \qquad
  s^{(i)} = s\!\left(\cdot,\,\cdot\,;\,u^{(i)}_0,\kappa^{(i)}\right)
  \in L^2(\Omega\times\mathcal{T}),
  \label{eq:dataset}
\end{equation}
where $u_0^{(i)} \in \mathcal{U} = L^2(\Omega)$ is the initial condition and $\kappa^{(i)} \in \cK$ the physical parameter.
Each trajectory $s^{(i)}$ is discretized on $N_y = N_t N_x$ quadrature points
$\{y_j = (x_j, t_j),\, w_j\}_{j=1}^{N_y}$ covering the product grid
$\Omega \times \mathcal{T}$, with norm $\normsq{f} = \sum_j w_j f(y_j)^2$.
We use the term \emph{trajectory} to refer to any solution field indexed by $(u_0,\kappa)$;
the formulation applies equally to purely spatial problems (where $\mathcal{T}$ collapses to a single point)
or to any other product structure on the output domain.


The dataset $\cD$ captures $\cG$ through a finite collection of trajectories; the underlying operator is defined as follows.

\begin{definition}[Solution operator]
\label{def:operator}
The \emph{solution operator}
 \[ \cG: \mathcal{U} \times \cK \to L^2(\Omega \times \T) \]
maps each pair $(u_0,\kappa)$ to the full space-time solution:
$\cG(u_0,\kappa) = s(\,\cdot\,,\,\cdot\,;\,u_0,\kappa)$.
\end{definition}


The operator $\cG$ is unknown and expensive to evaluate directly. The \emph{operator learning problem} is to construct $\hat{\cG} \approx \cG$ from $\cD$ alone, generalizing to unseen $(u_0,\kappa) \notin \cD$ without solving the PDE.





\subsection*{The multi-regime assumption}

When solutions change character fundamentally across $\cK$, the optimal low-dimensional representation differs from regime to regime, and no single global basis captures the full solution diversity. We formalize this observation via a latent variable model over the trajectory ensemble.

\begin{definition}[Generative model]
\label{def:generative}
Each trajectory arises from one of $M$ latent regimes. The regime label $z_i \in \{1,\ldots,M\}$ is unobserved and drawn from a categorical prior,
\begin{equation}
  z_i \sim \mathrm{Categorical}(\pi_1,\ldots,\pi_M),
  \qquad \pi_m \geq 0,\quad \textstyle\sum_m \pi_m = 1.
  \label{eq:prior}
\end{equation}
Conditioned on $z_i = m$, the trajectory is distributed as
\begin{equation}
 ( s^{(i)} \mid z_i{=}m) \;\sim\;
  \cN\!\bigl(\hats_m^{(i)},\;\sigma^2 W^{-1}\bigr),
  \label{eq:likelihood}
\end{equation}
where $W = \mathrm{diag}(w_1,\ldots,w_{N_y})$ is the quadrature weight matrix, $\sigma^2 > 0$ controls the softness of regime assignments, and $\hats_m^{(i)}$ is the regime-$m$ approximation to $s^{(i)}$.
\end{definition}


The generative model leaves $\hats_m^{(i)}$ unspecified. We require only that it admits a \emph{low-rank trajectory decomposition}:
\begin{equation}
  \hats_m^{(i)}(y)
  \;=\; \phi_0^{(m)}(y)
  + \sum_{p=1}^{P_m}
    \lambda_p^{(m,i)}\,\phi_p^{(m)}(y),
  \qquad y = (x,t) \in \Omega\times\mathcal{T},
  \label{eq:separable}
\end{equation}
where $\phi_p^{(m)} : \Omega\times\mathcal{T}\to\mathbb{R}$ are spatiotemporal
mode functions and $\lambda_p^{(m,i)} \in \mathbb{R}$ a scalar amplitude
for trajectory~$i$ on mode~$p$.
This form encompasses POD~\citep{Lu2022} and learned neural bases~\citep{Mou2026};
in this work we adopt the latter.

\begin{definition}[Regime approximator]
\label{def:approx}
The parameters of regime $m$ are
$\Phi_m = \{\phi_p^{(m)}\}_{p=0}^{P_m}$
and scalar amplitudes
$\Lambda_m = \{\lambda_p^{(m,i)}\}_{p=1,\,i=1}^{P_m,\,N}$,
where $\phi_0^{(m)} : \Omega\times\mathcal{T}\to\mathbb{R}$ is the regime mean field,
$\{\phi_p^{(m)}\}_{p \geq 1}$ spatiotemporal mode functions,
$\lambda_p^{(m,i)}\in\mathbb{R}$ the scalar amplitude of trajectory~$i$ on mode~$p$,
and $P_m \leq P_{\max}$ the adaptive mode count.
\end{definition}

The learning problem decomposes into two phases: discovering the latent regime structure from $\cD$ \emph{offline}, then training a deployable operator for fast inference \emph{online}.


\subsection*{Phase 1: Offline regime discovery}

Since the regime labels $z_i$ in~\eqref{eq:likelihood} are unobserved,
we recover the bases $\{\Phi_m\}$, mixture weights $\pi$, and noise
level $\sigma^2$ by maximizing the observed-data log-likelihood:
\begin{equation}
  (\Phi^*,\pi^*,\sigma^{2*})
  \;=\; \operatorname*{arg\,max}_{\Phi,\,\pi,\,\sigma^2}
  \;\sum_{i=1}^N \log \sum_{m=1}^M
    \frac{\pi_m\,\det(W)^{1/2}}{(2\pi\sigma^2)^{N_y/2}}
    \exp\!\!\left(-\frac{\normsq{s^{(i)}-\hats_m^{(i)}}}{2\sigma^2}\right).
  \label{eq:mle}
\end{equation}
The solution yields soft regime assignments
$\gamma_{im}^* = p(z_i{=}m \mid s^{(i)},\Phi^*,\pi^*,\sigma^{2*})$
that encode the discovered regime structure of $\cD$
and serve as supervision for Phase~2.



\subsection*{Phase 2: Online operator learning}

With bases $\{\Phi_m^*\}$ and responsibilities $\{\gamma_{im}^*\}$
frozen, we train a \emph{GatingNet} and $M$ \emph{BranchNets} that
map a new input $(u_0, \kappa)$ to a weighted combination of local
regime predictions. The training objective balances three goals:
reconstruction accuracy on $\cD$, consistency of the learned gating
with the offline regime structure $\{\gamma_{im}^*\}$, and diversity
across regimes to prevent collapse. The precise architecture and loss
are given in Section~\ref{sec:theory}.

The quality of the final approximation is measured by the mean
relative $L^2$ error on held-out test snapshots:
\begin{equation}
  \varepsilon \;=\; \frac{1}{N_\mathrm{test}}
  \sum_{i=1}^{N_\mathrm{test}}
  \frac{\norm{s^{(i)} - \hat{s}(\cdot;\,u_0^{(i)},\kappa^{(i)})}}
       {\norm{s^{(i)}}} .
  \label{eq:metric}
\end{equation}

\textit{The algorithm} solves Phase~1 via EM and 
Phase~2 via gradient descent; details are in Section~\ref{sec:theory}.
\newpage















\section{Theory}\label{sec:theory}



All methods studied in this work share a common two-phase structure:
an \emph{offline} basis-fitting phase that constructs a low-rank
approximation to the solution manifold from snapshots in $\mathcal{D}$,
and an \emph{online} operator-learning phase that trains a branch
network to predict expansion coefficients from a new input
$(u_0, \kappa)$.
The four methods differ along two independent axes:

\begin{itemize}
  \item \textbf{Basis type.}
    A \emph{linear} (POD) basis uses SVD of the (weighted) snapshot
    matrix; a \emph{nonlinear} (NeuralPOD) basis learns spatiotemporal
    mode networks by sequential residual minimization.
  \item \textbf{Number of regimes.}
    A \emph{single-regime} operator fits one global basis to all
    snapshots; a \emph{multi-regime} operator (TEMPO) discovers $M$
    latent regimes via EM and assigns a dedicated basis to each.
\end{itemize}

\noindent
This gives four combinations:
POD-DeepONet~\cite{Lu2022} (linear, global),
NeuralPOD-DeepONet (nonlinear, global, our extension of~\cite{Mou2026}),
TEMPO(POD) (linear, $M$ regimes), and
TEMPO(NeuralPOD) (nonlinear, $M$ regimes).
Sections~\ref{sec:bases}--\ref{sec:online} develop each component.


\subsection{Basis Representations}\label{sec:bases}


\subsubsection{POD Basis}

Let $\bar{s} = \tfrac{1}{N}\sum_i s^{(i)}$ and let
$\tilde{S} \in \mathbb{R}^{N_y \times N}$ have columns
$\tilde{s}^{(i)} = s^{(i)} - \bar{s}$.
The rank-$K$ truncated SVD $\tilde{S} \approx U_K \Sigma_K V_K^\top$
yields orthonormal modes $\{\phi_k\}_{k=1}^K$ (columns of $U_K$) and
coefficients $\lambda_k^{(i)} = [\Sigma_K V_K^\top]_{ki}$; the
approximation
\begin{equation}
  \hat{s}^{(i)} = \bar{s} + \sum_{k=1}^{K} \lambda_k^{(i)}\,\phi_k
  \label{eq:pod_approx}
\end{equation}
is optimal in the $L^2$ sense among all rank-$K$ affine representations.

For TEMPO(POD), the M-step~\eqref{eq:mstep_phi} is solved analytically.
Given responsibilities $\{\gamma_{im}\}$, replace $\bar{s}$ with the
weighted mean $\bar{s}_m = \sum_i \gamma_{im} s^{(i)}$ and apply SVD to
the weighted centered matrix with columns
$\sqrt{\gamma_{im}}\,(s^{(i)}-\bar{s}_m)$.
This is the closed-form linear analogue of
Algorithm~\ref{alg:npod}.

\subsubsection{NeuralPOD Basis}\label{sec:neuralpod_basis}

To overcome the linearity constraint, we parametrize each mode
function $\phi_k : \Omega{\times}\mathcal{T} \to \mathbb{R}$ as a
\emph{Fourier feature network}~\cite{Tancik2020}:
\begin{equation}
  \phi_k(y) \;=\; f_{\theta_k}\!\bigl(\psi(y)\bigr),
  \qquad
  \psi(y) \;=\;
  \bigl[\sin(2\pi By),\;\cos(2\pi By)\bigr],
  \label{eq:ffn}
\end{equation}
where $B \in \mathbb{R}^{F \times d}$ is a fixed random matrix drawn
once at initialization and $f_{\theta_k} : \mathbb{R}^{2F} \to \mathbb{R}$
is a small MLP with Tanh activations.
Plain MLPs exhibit spectral bias~\cite{Rahaman2019} and cannot capture
high-frequency features such as shock fronts; the Fourier encoding
resolves this.
For multi-scale coverage, the rows of $B$ are partitioned evenly across
$L$ scales $s_1 < \cdots < s_L$: the $\ell$-th block is drawn from
$\mathcal{N}(0,\,s_\ell^2 I)$.
The mean field $\phi_0$ uses the same architecture.

The scalar amplitudes $\lambda_k^{(i)} \in \mathbb{R}$ are direct
learnable parameters, one per training trajectory, optimized jointly
with $\phi_k$.
After Phase~1 they are frozen and serve as regression targets for the
BranchNet in Phase~2.

Modes are learned by \emph{greedy residual minimization}~\cite{Mou2026}.
Let $r^{(0,i)} = s^{(i)} - \phi_0$, where $\phi_0$ is fitted to the snapshot mean $\bar{s}$ defined above.
For each subsequent mode $k \ge 1$, given the residual
\begin{equation}
  r^{(k-1,i)}(y)
  = s^{(i)}(y)
  - \phi_0(y)
  - \sum_{j=1}^{k-1}\lambda_j^{(i)}\,\phi_j(y),
  \label{eq:residual_global}
\end{equation}
we solve
\begin{equation}
  \min_{\phi_k,\,\{\lambda_k^{(i)}\}}
  \sum_{i=1}^{N} \sum_j w_j
  \Bigl(\lambda_k^{(i)}\,\phi_k(y_j) - r_j^{(k-1,i)}\Bigr)^{\!2}
  \label{eq:neuralpod_loss}
\end{equation}
(the global case with uniform weights; the $\gamma$-weighted
generalisation used in TEMPO is Algorithm~\ref{alg:npod}).
Mode addition stops when the residual falls below fraction $\tau$ of
the initial variance,
\begin{equation}
  \sum_{i=1}^N \|r^{(k,i)}\|_w^2
  \;<\;
  \tau \cdot \sum_{i=1}^N \|s^{(i)} - \bar{s}\|_w^2,
  \label{eq:npod_tol}
\end{equation}
or $k = K_{\max}$.
Each mode captures the largest remaining component of variance without
constraining the modes to a fixed linear subspace, making this the
nonlinear analogue of Gram--Schmidt orthogonalization in the sense of
sequential variance extraction.

\subsection{Single-Regime Operators}\label{sec:single}


\subsubsection{POD-DeepONet}

POD-DeepONet~\cite{Lu2022} replaces the learned trunk network of
DeepONet~\cite{Lu2019} with the fixed POD basis.
After computing the global modes $\{\phi_k\}_{k=1}^K$ and training
coefficients $\{\lambda^{(i)}\}$ offline, a branch network
$\mathcal{B}_\theta : \mathbb{R}^{m_s + d_\kappa} \to \mathbb{R}^K$
is trained online by regression against these coefficients:
\begin{equation}
  \mathcal{L}_{\mathrm{branch}}
  = \frac{1}{N}\sum_{i=1}^{N}
    \bigl\|\mathcal{B}_\theta(u_0^{(i)},\,\kappa^{(i)})
    - \lambda^{(i)}\bigr\|^2.
  \label{eq:pod_branch_loss}
\end{equation}
The prediction at any $(x,t) \in \Omega \times \mathcal{T}$ is
\begin{equation}
  \hat{s}(x,t;\,u_0,\kappa)
  = \phi_0(x,t)
  + \sum_{k=1}^{K}
    \mathcal{B}_{\theta,k}(u_0,\kappa)\,\phi_k(x,t).
  \label{eq:pod_predict}
\end{equation}

\subsubsection{NeuralPOD-DeepONet}\label{sec:neuralpod_deeponet}

Mou et al.~\cite{Mou2026} introduce NeuralPOD as a dimensionality
reduction for PDE solutions, establishing the greedy basis-fitting
procedure of Section~\ref{sec:neuralpod_basis}.
However, they do not construct a deployable solution operator: no
network is proposed to predict the mode coefficients for an unseen
initial condition.
We close this gap by substituting the NeuralPOD basis into the
POD-DeepONet framework.

\paragraph{Phase 1 (offline).}
Train the global NeuralPOD basis
$\{\phi_0, \phi_1, \ldots, \phi_K\}$
on all $N$ training snapshots by the procedure of
Section~\ref{sec:neuralpod_basis}, obtaining coefficients
$\boldsymbol{\lambda}^{(i)} \in \mathbb{R}^K$ for each snapshot.

\paragraph{Phase 2 (online).}
Fix the NeuralPOD basis. Train a branch network
$\mathcal{B}_\theta : \mathbb{R}^{m_s + d_\kappa} \to \mathbb{R}^K$
by minimizing~\eqref{eq:pod_branch_loss} with the NeuralPOD
coefficients as regression targets.

\paragraph{Prediction.}
The inference formula is identical to~\eqref{eq:pod_predict}, with
the NeuralPOD networks substituted for the POD modes.
Because the nonlinear basis captures solution structure that lies
off any fixed linear subspace, this operator spans a strictly richer
function space than POD-DeepONet at the same mode count $K$.


\subsection{TEMPO: Offline Regime Discovery}\label{sec:offline}


A single global basis — whether linear or nonlinear — must compromise
between qualitatively distinct solution regimes.
For problems such as Burgers' equation, where small viscosity produces
sharp shock structures and large viscosity yields smooth diffusion, the
optimal basis geometry differs fundamentally between regimes.
TEMPO avoids this compromise by constructing a dedicated basis for each
of $M$ latent regimes, discovered from unlabelled trajectories via
Expectation--Maximization on the generative model of
Definition~\ref{def:generative}.

Under regime $m$, the approximation follows~\eqref{eq:separable}:
\begin{equation}
  \hat{s}_m^{(i)}(y)
  = \phi_0^{(m)}(y)
  + \sum_{p=1}^{P_m}
    \lambda_p^{(m,i)}\,\phi_p^{(m)}(y),
  \qquad y\in\Omega\times\mathcal{T}.
  \label{eq:regime_approx}
\end{equation}
Each trajectory $s^{(i)}\in\mathbb{R}^{N_y}$ is treated as a single data point on
the product grid $\{y_j=(x_j,t_j)\}$, so the temporal structure is encoded in the
coordinate $y$ rather than separated as a per-snapshot argument.
The scalar amplitudes $\{\lambda_p^{(m,i)}\}$ are learned jointly with the mode
functions $\{\phi_p^{(m)}\}$ during offline training and serve as regression targets
for the online BranchNet $\mathcal{B}_m(u_0,\kappa)\to\mathbb{R}^{P_m}$,
which predicts these amplitudes for unseen inputs at inference time.

\paragraph{Initialization.}
A global POD is computed on all $N$ trajectories; each trajectory is
projected to a $P$-dimensional coefficient vector
$\alpha^{(i)} \in \mathbb{R}^P$.
Running a standard GMM on $\{\alpha^{(i)}\}$ yields initial
responsibilities $\gamma_{im}^{(0)}$, weights $\pi_m^{(0)}$,
and per-regime statistics $(\mu_m^{(0)}, \Sigma_m^{(0)})$.
Each regime basis $\Phi_m^{(0)}$ is then fitted to the
responsibility-weighted ensemble: for TEMPO(NeuralPOD) via
\textsc{WeightedNeuralPOD} (Algorithm~\ref{alg:npod}); for TEMPO(POD)
via the weighted SVD of Section~\ref{sec:bases}.
The noise level $\sigma^2$ is calibrated from the initial
reconstructions rather than treated as a free hyperparameter:
\begin{equation}
  \sigma^2 \;\leftarrow\;
  \frac{1}{2N}\sum_{i=1}^N\sum_{m=1}^M
  \gamma_{im}^{(0)}\,\bigl\|s^{(i)} - \hat{s}_m^{(i)}\bigr\|_w^2,
  \label{eq:sigma2_calib}
\end{equation}
which ties the softness of the E-step directly to the quality of the
initial bases.

\paragraph{E-step.}
At iteration $q$, the posterior regime responsibility for trajectory $i$
follows directly from Bayes' theorem applied to the Gaussian likelihood
of Definition~\ref{def:generative}:
\begin{equation}
  \gamma_{im}^{(q)}
  = \frac{
    \pi_m^{(q-1)}\,
    \exp\!\Bigl(-\dfrac{\|s^{(i)}-\hat{s}_m^{(i)}\|_w^2}{2\sigma^2}\Bigr)
  }{
    \sum_{j=1}^{M}
    \pi_j^{(q-1)}\,
    \exp\!\Bigl(-\dfrac{\|s^{(i)}-\hat{s}_j^{(i)}\|_w^2}{2\sigma^2}\Bigr)
  },
  \label{eq:estep}
\end{equation}
where $\hat{s}_m^{(i)}$ is evaluated via~\eqref{eq:regime_approx}
using $\Phi_m^{(q-1)}$ and $\Lambda_m^{(q-1)}$.
Crucially, as the bases improve in the M-step, the
reconstructions $\hat{s}_m^{(i)}$ sharpen and the E-step
automatically refines the regime assignments in response.
This coupling distinguishes TEMPO from a GMM applied to a fixed
feature space.

\begin{remark}[Heteroscedastic extension]
When $\kappa$ drives order-of-magnitude variation in solution amplitude,
the homoscedastic likelihood~\eqref{eq:likelihood} is poorly calibrated:
regimes with larger solutions dominate the E-step.
Replacing the absolute error $\|s^{(i)}-\hat{s}_m^{(i)}\|_w^2$ with the
relative error $\|s^{(i)}-\hat{s}_m^{(i)}\|_w^2\,/\,\|s^{(i)}\|_w^2$
in both the E-step and $\sigma^2$ calibration yields a scale-invariant model.
When $\kappa$ is observed, the GMM initialization can be replaced by soft
assignments in $\log\kappa$ space; an optional $\kappa$-prior term
(weighted by $\lambda \ge 0$) in the E-step anchors regime assignments
to known parameter structure throughout EM.
\end{remark}

\paragraph{M-step.}
Maximizing the expected complete-data log-likelihood with
responsibilities fixed yields a closed-form update for the mixture
weights,
\begin{equation}
  \pi_m^{(q)} \leftarrow \frac{N_m^{(q)}}{N},
  \qquad N_m^{(q)} = \sum_{i=1}^{N}\gamma_{im}^{(q)},
  \label{eq:mstep_pi}
\end{equation}
and a weighted reconstruction problem for each basis:
\begin{equation}
  \min_{\Phi_m,\,\Lambda_m}
  \sum_{i=1}^{N} \gamma_{im}^{(q)}\,\|s^{(i)} - \hat{s}_m^{(i)}\|_w^2,
  \label{eq:mstep_phi}
\end{equation}
solved by weighted SVD for TEMPO(POD) or
\textsc{WeightedNeuralPOD} (Algorithm~\ref{alg:npod}) for
TEMPO(NeuralPOD).
Because the NeuralPOD subproblem is minimized by gradient descent
rather than to global optimality, TEMPO is a Generalized EM
algorithm~\cite{Neal1998}: it guarantees monotone non-decrease of the
observed-data log-likelihood at every iteration under the
sufficient condition that each M-step does not decrease the
Q-function.

\paragraph{Adaptive basis update.}
Retraining all $M$ NeuralPOD bases at every EM iteration is
computationally prohibitive.
We introduce a distribution-shift criterion to determine whether each
basis requires updating.
Let $\alpha^{(i)}$ denote the global POD projection of $s^{(i)}$,
computed once during initialization and fixed thereafter.
The responsibility-weighted first and second moments of regime $m$ are
\begin{align}
  \mu_m^{(q)}
  &= \frac{1}{N_m^{(q)}}\sum_{i=1}^N
    \gamma_{im}^{(q)}\,\alpha^{(i)},
  \label{eq:mu}\\
  \Sigma_m^{(q)}
  &= \frac{1}{N_m^{(q)}}\sum_{i=1}^N
    \gamma_{im}^{(q)}\,
    (\alpha^{(i)}-\mu_m^{(q)})(\alpha^{(i)}-\mu_m^{(q)})^\top.
  \label{eq:Sigma}
\end{align}
The distribution-shift criterion measures the relative change in these
moments between consecutive iterations:
\begin{equation}
  \Delta_m^{(q)}
  = \frac{\|\mu_m^{(q)} - \mu_m^{(q-1)}\|_2}
         {\|\mu_m^{(q-1)}\|_2 + \epsilon}
  + \frac{\|\Sigma_m^{(q)} - \Sigma_m^{(q-1)}\|_F}
         {\|\Sigma_m^{(q-1)}\|_F + \epsilon}.
  \label{eq:Delta}
\end{equation}
Note that $\Delta_m^{(q)}$ tracks the shift in the \emph{data content}
of regime $m$ — the trajectories it effectively owns — not merely
the volume of responsibility reassignments.
Two trajectories with very different $\alpha^{(i)}$ swapping regime
membership produce a large $\Delta_m^{(q)}$, triggering a basis
update; two similar snapshots doing the same do not.
Given thresholds $0 < \varepsilon_s < \varepsilon_l$, the update
decision at iteration $q$ follows a three-way rule:

\begin{center}
\renewcommand{\arraystretch}{1.3}\small
\begin{tabular}{lll}
\toprule
Condition & Action & Rationale \\
\midrule
$\Delta_m^{(q)} < \varepsilon_s$
  & \textsc{Skip}
  & Data manifold unchanged; basis remains accurate \\
$\varepsilon_s \le \Delta_m^{(q)} < \varepsilon_l$
  & \textsc{Incremental}
  & Moderate drift; add or prune one mode \\
$\Delta_m^{(q)} \ge \varepsilon_l$
  & \textsc{Full rerun}
  & Large shift; rebuild from random initialization \\
\bottomrule
\end{tabular}
\end{center}

\noindent
In the \textsc{Incremental} case, we compute the residual of the
current basis under the updated responsibilities $\gamma^{(q)}$, add
a new mode if $\sum_i \gamma_{im}^{(q)}\|r_m^{(P_m,i)}\|_w^2 \ge \tau\cdot\sum_i \gamma_{im}^{(q)}\|s^{(i)}-\bar{s}_m\|_w^2$,
and prune the last mode if its weighted contribution falls below
$\varepsilon_p$.
The \textsc{Full rerun} case uses a cold-start random initialization.
When $\Delta_m^{(q)} \ge \varepsilon_l$, the parameters
$\Phi_m^{(q-1)}$ are a local minimum of the \emph{previous}
weighted objective and may be in a basin that is entirely
misaligned with the new weighted objective; warm-starting would
trap gradient descent there, encoding the old regime's features
in the new basis.

\paragraph{Convergence and model selection.}
EM terminates when $\max_m \Delta_m^{(q)} < \varepsilon_c$.
The number of regimes $M \in \{2,3,4\}$ is selected by jointly
minimizing BIC on the GMM fit and mean assignment entropy;
see Table~\ref{tab:regime_selection}.

\begin{algorithm}[t]
\caption{\textsc{WeightedNeuralPOD}$\bigl(
  \{s^{(i)}, \gamma_{im}\}_{i=1}^N,\;
  \tau\bigr)$}
\label{alg:npod}
\begin{algorithmic}[1]
\State $\bar{s}^{(m)} \leftarrow \displaystyle\sum_{i=1}^N \gamma_{im}\,s^{(i)}$
  \hfill\Comment{responsibility-weighted mean trajectory}
\State Train $\phi_0^{(m)}$ by minimizing
  $\displaystyle\sum_j w_j\bigl(\phi_0^{(m)}(y_j) - \bar{s}^{(m)}_j\bigr)^2$
\State $r^{(0,i)} \leftarrow s^{(i)} - \phi_0^{(m)}(y)$
  \enspace for all $i$;\enspace $p \leftarrow 0$
\State $V_m \leftarrow \displaystyle\sum_{i=1}^N \gamma_{im}\,\|s^{(i)}-\bar{s}^{(m)}\|_w^2$
  \hfill\Comment{initial unmodeled variance}
\While{$\displaystyle\sum_{i=1}^N \gamma_{im}\,\|r^{(p,i)}\|_w^2
  \;\ge\; \tau\, V_m$ \textbf{ and } $p < P_{\max}$}
  \State Jointly minimize over $\phi_{p+1}^{(m)}$ and
    $\{\lambda_{p+1}^{(m,i)}\}_{i=1}^N$:
    \[
      \sum_{i=1}^N \gamma_{im}
      \sum_j w_j\!\left(\lambda_{p+1}^{(m,i)}\,\phi_{p+1}^{(m)}(y_j)
      - r^{(p,i)}_j\right)^{\!2}
    \]
  \State $r^{(p+1,i)} \leftarrow r^{(p,i)}
    - \lambda_{p+1}^{(m,i)}\,\phi_{p+1}^{(m)}(y)$
    \enspace for all $i$
  \State $p \leftarrow p + 1$
\EndWhile
\State \Return $P_m \!\leftarrow\! p$,\enspace
  $\phi_0^{(m)}$,\enspace
  $\bigl\{\phi_p^{(m)},\,\{\lambda_p^{(m,i)}\}_{i=1}^N\bigr\}_{p=1}^{P_m}$
\end{algorithmic}
\end{algorithm}

\begin{algorithm}[t]
\caption{TEMPO --- Offline Phase}
\label{alg:tempo}
\begin{algorithmic}[1]
\Require $\mathcal{D} = \{u_0^{(i)}, \kappa^{(i)}, s^{(i)}\}_{i=1}^N$;\;
  $M,\;\varepsilon_s,\;\varepsilon_l,\;\varepsilon_p,\;\varepsilon_c,\;\tau$
\Ensure Regime bases $\{\Phi_m^*\}$,\;
  amplitudes $\{\Lambda_m^*\}$,\;
  responsibilities $\{\gamma_{im}^*\}$,\;
  weights $\{\pi_m^*\}$
%
\Statex \textit{Initialization}
\State Compute global POD on $\{s^{(i)}\}$;
  project to $\alpha^{(i)} \in \mathbb{R}^P$ for all $i$
\State Run GMM-EM on $\{\alpha^{(i)}\}$ $\to$
  $\gamma_{im}^{(0)},\;\pi_m^{(0)},\;\mu_m^{(0)},\;\Sigma_m^{(0)}$
\For{$m = 1,\ldots,M$}
  \State $\Phi_m^{(0)},\,\Lambda_m^{(0)} \leftarrow$
    \Call{WeightedNeuralPOD}
    {$\{s^{(i)},\gamma_{im}^{(0)}\},\,\tau$}
    \hfill\Comment{or weighted SVD for TEMPO(POD), see \S\ref{sec:bases}}
\EndFor
\State $\sigma^2 \leftarrow (2N)^{-1}\!\sum_{i,m}
  \gamma_{im}^{(0)}\,\|s^{(i)}{-}\hat{s}_m^{(i)}\|_w^2$
  \hfill\eqref{eq:sigma2_calib}
%
\Statex \textit{EM iterations}
\Repeat
  \Statex \quad\textit{E-step}
  \For{all $i = 1,\ldots,N$ and $m = 1,\ldots,M$}
    \State Compute $\hat{s}_m^{(i)}$ via~\eqref{eq:regime_approx}
      using $\Phi_m^{(q-1)}$ and $\Lambda_m^{(q-1)}$
    \State $\gamma_{im}^{(q)} \leftarrow$ Eq.~\eqref{eq:estep}
  \EndFor
  \Statex \quad\textit{M-step: mixture weights}
  \State $\pi_m^{(q)} \leftarrow N_m^{(q)}/N$
    for all $m$ \hfill\eqref{eq:mstep_pi}
  \Statex \quad\textit{M-step: bases}
  \For{$m = 1,\ldots,M$}
    \State Compute $\mu_m^{(q)}$~\eqref{eq:mu},\;
      $\Sigma_m^{(q)}$~\eqref{eq:Sigma},\;
      $\Delta_m^{(q)}$~\eqref{eq:Delta}
    \If{$\Delta_m^{(q)} < \varepsilon_s$}
      \State $\Phi_m^{(q)} \leftarrow \Phi_m^{(q-1)}$
        \hfill \textsc{Skip}
    \ElsIf{$\Delta_m^{(q)} < \varepsilon_l$}
      \hfill \textsc{Incremental}
      \State Compute $r_m^{(P_m)}$ under updated $\gamma^{(q)}$;\;
        $V_m^{(q)} \leftarrow \sum_i\gamma_{im}^{(q)}\|s^{(i)}{-}\bar{s}_m^{(q)}\|_w^2$
      \If{$\sum_i\gamma_{im}^{(q)}\|r_m^{(P_m,i)}\|_w^2 \ge \tau\cdot V_m^{(q)}$}
        \State Train mode $P_m{+}1$;\; $P_m \leftarrow P_m + 1$
      \EndIf
      \If{$\|\phi_{P_m}^{(m)}\|_w^2\cdot\displaystyle\sum_{i=1}^N
        \gamma_{im}^{(q)}\,\bigl(\lambda_{P_m}^{(m,i)}\bigr)^2 < \varepsilon_p$}
        \State $P_m \leftarrow P_m - 1$
      \EndIf
    \Else
      \hfill \textsc{Full rerun}
      \State $\Phi_m^{(q)},\,\Lambda_m^{(q)} \leftarrow$
        \Call{WeightedNeuralPOD}
        {$\{s^{(i)},\gamma_{im}^{(q)}\},\,\tau$}
    \EndIf
  \EndFor
\Until{$\max_m \Delta_m^{(q)} < \varepsilon_c$}
\State \Return $\Phi_m^* \!\leftarrow\! \Phi_m^{(q)}$,\;
  $\Lambda_m^* \!\leftarrow\! \Lambda_m^{(q)}$,\;
  $\gamma_{im}^* \!\leftarrow\! \gamma_{im}^{(q)}$,\;
  $\pi_m^* \!\leftarrow\! \pi_m^{(q)}$
\end{algorithmic}
\end{algorithm}


\subsection{TEMPO: Online Gated Operator}\label{sec:online}


After Algorithm~\ref{alg:tempo} converges, all regime bases
$\{\Phi_m^*\}$ and responsibilities $\{\gamma_{im}^*\}$ are frozen.
The online phase trains a \textbf{GatingNet}
$\mathcal{W}: \mathbb{R}^{m_s + d_\kappa} \to \Delta^{M-1}$
and $M$ \textbf{BranchNets}
$\mathcal{B}_m: \mathbb{R}^{m_s + d_\kappa} \to \mathbb{R}^{P_m}$,
each mapping sensor observations $u_0$ and parameter
$\kappa$ to regime weights and expansion coefficients, respectively.
The TEMPO prediction at query point $y = (x,t) \in \Omega\times\mathcal{T}$ is
\begin{equation}
  \hat{s}(y;\,u_0,\kappa)
  = \sum_{m=1}^{M} w_m(u_0,\kappa)
    \left[
      \phi_0^{(m)}\!(y)
      + \sum_{k=1}^{P_m}
        b_{m,k}(u_0,\kappa)\;
        \phi_k^{(m)}\!(y)
    \right],
  \label{eq:predict}
\end{equation}
where $\mathbf{w} = \mathcal{W}(u_0,\kappa)$ and
$\mathbf{b}_m = \mathcal{B}_m(u_0,\kappa)$.
The BranchNet outputs $\mathbf{b}_m$ are the online predictions of the offline
amplitudes $\boldsymbol{\lambda}_m^{(i)}$, enabling generalisation to unseen
$(u_0,\kappa)$.
The prediction is mesh-free and evaluable at any $y\in\Omega\times\mathcal{T}$
without retraining.

\paragraph{Training objective.}
Writing $w_m^{(i)} = w_m(u_0^{(i)},\kappa^{(i)})$, the online loss is
\begin{equation}
  \mathcal{L}_{\mathrm{online}}
  = \underbrace{
      \frac{1}{N}\sum_{i=1}^{N}
      \|s^{(i)} - \hat{s}(\cdot;\,u_0^{(i)},\kappa^{(i)})\|_w^2
    }_{\mathcal{L}_{\mathrm{data}}}
  +\;\lambda_{\mathrm{KL}}
    \underbrace{
      \frac{1}{N}\sum_{i,m}
      w_m^{(i)}\log\frac{w_m^{(i)}}{\gamma_{im}^*}
    }_{\mathcal{L}_{\mathrm{KL}}}
  -\;\lambda_{H}
    \underbrace{
      \frac{1}{N}\sum_{i,m}
      w_m^{(i)}\log w_m^{(i)}
    }_{-\mathcal{L}_{\mathrm{ent}}}.
  \label{eq:online_loss}
\end{equation}
$\mathcal{L}_{\mathrm{data}}$ drives reconstruction accuracy.
$\mathcal{L}_{\mathrm{KL}}$ anchors the online gating to the
offline regime partition: it penalizes deviation of $\mathbf{w}^{(i)}$
from the EM responsibilities $\boldsymbol{\gamma}_i^*$, so the
network cannot discard the discovered structure in pursuit of a
simpler routing.
$\mathcal{L}_{\mathrm{ent}}$ prevents collapse to a single expert,
ensuring all $M$ regime bases contribute.
The hyperparameters $\lambda_{\mathrm{KL}}, \lambda_H \ge 0$ trade
off these objectives.







\clearpage

\section{Computational Experiments}
\label{sec:exp}

We evaluate TEMPO on parametric PDE benchmarks, verifying three claims:
(i)~the offline EM phase recovers interpretable, physics-consistent regime
structure from unlabelled snapshots;
(ii)~locally adapted bases achieve lower reconstruction error than their
single-regime counterparts at equal mode count;
(iii)~TEMPO outperforms single-regime baselines in end-to-end operator accuracy.

\subsection{Experimental Setup}
\label{sec:setup}

\paragraph{Datasets.}
We use the PDEBench benchmark~\cite{Takamoto2022} across three problems.

\textbf{1D Burgers' equation.}
\begin{equation}
  u_t + u\,u_x = \nu\,u_{xx},
  \qquad x\in(-1,1),\quad t\in(0,2],
  \label{eq:burgers}
\end{equation}
with periodic boundary conditions and random initial conditions drawn from a
fixed distribution.
Viscosity $\nu\in\{0.001,\,0.1,\,1.0\}$ spans qualitatively distinct
solution regimes: smooth diffusive profiles at $\nu{=}1.0$ and
near-discontinuous shock structures at $\nu{=}0.001$.
Each value provides $N{=}9{,}500$ trajectories discretised on $N_x{=}1{,}024$
spatial points and $N_t{=}201$ time steps, yielding $28{,}500$ trajectories in total.

\textbf{2D Darcy flow.}
\begin{equation}
  -\nabla\cdot\bigl(a(x)\,\nabla u(x)\bigr) = \beta,
  \qquad x\in(0,1)^2,
  \label{eq:darcy}
\end{equation}
with zero Dirichlet boundary conditions.
Here $a(x)$ is a spatially varying permeability field sampled from a
random sum of Gaussians, and $\beta\in\{0.01,\,0.1,\,1.0,\,10.0,\,100.0\}$
is the constant source term.
Each value provides $N{=}8{,}000$ training and $1{,}000$ test samples
discretised on a $128{\times}128$ uniform grid ($N_y{=}16{,}384$ spatial points),
$40{,}000$ trajectories in total.
The operator maps $a(x)$ to the pressure solution $u(x)$.

\paragraph{Methods.}
We compare five methods, all implemented and trained by us on the same data
splits.
DeepONet~\cite{Lu2019} and POD-DeepONet~\cite{Lu2022} serve as single-regime
linear baselines.
NeuralPOD-DeepONet (\textit{ours}) extends POD-DeepONet with the Fourier
NeuralPOD basis of Section~\ref{sec:bases}.
TEMPO is evaluated in two variants:
\textbf{TEMPO(POD)} uses responsibility-weighted SVD as the regime basis, and
\textbf{TEMPO(NeuralPOD)} uses Fourier NeuralPOD bases, capturing nonlinear
and parameter-dependent solution structure.

\paragraph{Metric.}
All methods are evaluated by the mean relative $L^2$ error $\varepsilon$
defined in~\eqref{eq:metric}, computed per viscosity value and averaged
over $\nu\in\{0.001,\,0.01,\,0.1,\,1.0\}$.
We additionally report RMSE and MSE.

%==========================================================================
\subsection{Single-Regime Operators}
\label{sec:single_exp}
%==========================================================================

We characterise the single-regime baselines on the 1D Burgers' dataset,
separating basis quality (Phase~1) from operator accuracy (Phase~2).
This decomposition isolates where reconstruction error originates and
provides ablation points for the multi-regime analysis in
Section~\ref{sec:tempo_exp}.

\subsubsection{Phase~1: Basis Quality}
\label{sec:phase1_exp}

\paragraph{POD.}
Figure~\ref{fig:pod_phase1} shows the singular value spectrum and cumulative
explained variance.
The solution energy is heavily concentrated: $K{=}30$ modes capture $99.99\%$
of the total variance, achieving a Phase~1 relative $L^2$ reconstruction
error of $1.47\%$.
This optimal linear compression serves as an upper bound for what any
linear single-regime basis can achieve.

\begin{figure}[h]
\centering
\includegraphics[width=\textwidth]{TEMPO/main/pictures/pod_phase1.png}
\caption{POD basis on 1D Burgers' ($\nu{=}1.0$, $K{=}30$ modes).
  \textit{Left:} singular value spectrum.
  \textit{Right:} cumulative explained variance; $K{=}30$ modes attain
  $99.99\%$.}
\label{fig:pod_phase1}
\end{figure}

\paragraph{NeuralPOD.}
The steep initial decay confirms that the Fourier NeuralPOD basis concentrates
variance efficiently; subsequent modes capture incrementally finer solution
structure until the relative residual~\eqref{eq:npod_tol} falls below $\tau$.
Figure~\ref{fig:basis_rec} shows the learned spatial modes alongside
representative reconstructions.
Figure~\ref{fig:npod_phase1} shows the weighted residual norm as a function of
greedy mode count.


\begin{figure}[h]
\centering
\includegraphics[width=\textwidth]{TEMPO/main/pictures/Neural_POD_rec_basis.png}
\caption{Learned Fourier NeuralPOD basis functions and representative
  reconstructions on 1D Burgers' ($\nu{=}1.0$).}
\label{fig:basis_rec}
\end{figure}

\begin{figure}[h]
\centering
\includegraphics[width=\textwidth]{TEMPO/main/pictures/npod_phase1.png}
\caption{Fourier NeuralPOD basis on 1D Burgers' ($\nu{=}1.0$):
  weighted residual norm as a function of greedy mode count.}
\label{fig:npod_phase1}
\end{figure}

\subsubsection{Phase~2: Operator Learning}
\label{sec:phase2_exp}

With bases frozen, branch networks are trained by
minimising~\eqref{eq:pod_branch_loss}.
Figures~\ref{fig:pod_phase2_loss} and~\ref{fig:npod_phase2_loss} show the
MSE training curves for POD-DeepONet and NeuralPOD-DeepONet, respectively.
Table~\ref{tab:results} (rows 2--3) reports test-set accuracy;
Figures~\ref{fig:pod_err_dist} and~\ref{fig:npod_err_dist} show the full
test-set error distributions.
Figures~\ref{fig:pod_reconstruction} and~\ref{fig:npod_reconstruction} show
representative space--time predictions alongside point-wise absolute errors
for three randomly selected test snapshots.

\begin{figure}[h]
\centering
\includegraphics[width=\textwidth]{TEMPO/main/pictures/pod_phase2_loss.png}
\caption{Phase~2 MSE training curve for POD-DeepONet on 1D Burgers'
  ($\nu{=}1.0$).}
\label{fig:pod_phase2_loss}
\end{figure}

\begin{figure}[h]
\centering
\includegraphics[width=\textwidth]{TEMPO/main/pictures/npod_phase2_loss.png}
\caption{Phase~2 MSE training curve for NeuralPOD-DeepONet on 1D Burgers'
  ($\nu{=}1.0$).}
\label{fig:npod_phase2_loss}
\end{figure}

\begin{figure}[h]
\centering
\includegraphics[width=\textwidth]{TEMPO/main/pictures/pod_err_dist.png}
\caption{POD-DeepONet test-set error distribution on 1D Burgers' ($\nu{=}1.0$).
  \textit{Left:} histogram with mean (red) and median (blue).
  \textit{Right:} sorted relative $L^2$ errors.}
\label{fig:pod_err_dist}
\end{figure}

\begin{figure}[h]
\centering
\includegraphics[width=\textwidth]{TEMPO/main/pictures/npod_err_dist.png}
\caption{NeuralPOD-DeepONet test-set error distribution on 1D Burgers' ($\nu{=}1.0$).
  Layout as in Figure~\ref{fig:pod_err_dist}.}
\label{fig:npod_err_dist}
\end{figure}

\begin{figure}[h]
\centering
\includegraphics[width=\textwidth]{TEMPO/main/pictures/pod_reconstruction.png}
\caption{POD-DeepONet: representative reconstructions on 1D Burgers'
  ($\nu{=}1.0$). Each row shows one test snapshot: ground truth
  (\textit{left}), prediction (\textit{centre}), and point-wise absolute
  error (\textit{right}).}
\label{fig:pod_reconstruction}
\end{figure}

\begin{figure}[h]
\centering
\includegraphics[width=\textwidth]{TEMPO/main/pictures/npod_reconstruction.png}
\caption{NeuralPOD-DeepONet: representative reconstructions on 1D Burgers'
  ($\nu{=}1.0$). Layout as in Figure~\ref{fig:pod_reconstruction}.}
\label{fig:npod_reconstruction}
\end{figure}

%==========================================================================
\subsection{TEMPO: Multi-Regime Operator}
\label{sec:tempo_exp}
%==========================================================================

A single basis — linear or nonlinear — must allocate its capacity across
all solution regimes simultaneously.
On 1D Burgers' with $\nu\in\{0.001,\ldots,1.0\}$, this forces a compromise:
modes well-adapted to smooth diffusive profiles are poorly suited to
near-discontinuous shock structures.
TEMPO avoids this by discovering the latent regime partition offline and
constructing a dedicated basis for each regime.

\subsubsection{Offline Phase: Regime Discovery}
\label{sec:offline_exp}

\paragraph{Number of regimes.}
We select $M$ from $\{2,3,4\}$ by jointly minimising BIC on the EM-GMM fit
and mean assignment entropy
$H = -\frac{1}{N}\sum_{i,m}\gamma_{im}\log\gamma_{im}$,
where lower $H$ indicates more confident regime assignments.
Table~\ref{tab:regime_selection} reports both criteria;
$M{=}3$ achieves the best trade-off and is used throughout.

\begin{table}[h]
\centering
\caption{Regime selection on 1D Burgers'.}
\label{tab:regime_selection}
\begin{tabular}{ccc}
\toprule
$M$ & BIC & $H$ \\
\midrule
2          & -- & -- \\
\textbf{3} & $\mathbf{4.60{\times}10^{4}}$ & \textbf{0.908} \\
4          & -- & -- \\
\bottomrule
\end{tabular}
\end{table}

\paragraph{Regime structure.}
Figure~\ref{fig:umap_em} shows the UMAP embedding and the GMM
fit in POD coefficient space $(\alpha_1,\alpha_2)$.
Crucially, the three discovered regimes do \emph{not} coincide with
individual viscosity values: each manifold branch spans multiple~$\nu$,
confirming that the EM partition reflects intrinsic solution geometry rather
than the physical parameter directly.

After 10 EM iterations the mixture weights are
$\boldsymbol{\pi}{=}(0.429,\,0.336,\,0.235)$, indicating a balanced,
non-degenerate partition.
Regimes~1 and~2 have converged ($\Delta_m < 0.013$);
regime~3 continues to adapt ($\Delta_3{=}0.310$) and EM is ongoing.

\begin{figure}[h]
\centering
\includegraphics[width=\textwidth]{TEMPO/main/pictures/TEMPO_phase1.png}
\caption{TEMPO regime structure ($M{=}3$).
  \textit{Left:} UMAP coloured by regime assignment.
  \textit{Centre:} same embedding coloured by~$\nu$; each branch spans
  multiple viscosity values.
  \textit{Right:} GMM fit in POD coefficient space $(\alpha_1,\alpha_2)$;
  stars mark centroids~$\mu_m$.}
\label{fig:umap_em}
\end{figure}

\paragraph{Basis quality per regime.}
Figure~\ref{fig:training_dynamics} shows the Fourier NeuralPOD training
dynamics for a representative regime: the residual MSE decays rapidly,
with modes~1--2 capturing the dominant variance and diminishing contributions
beyond mode~4.
Figure~\ref{fig:basis_rec} shows the learned spatial modes alongside
representative reconstructions; modes are smooth and globally supported for
large~$\nu$ and sharply localised near the shock front for small~$\nu$,
demonstrating the parameter-adaptivity of the regime-specific NeuralPOD basis.



\clearpage
\subsubsection{Online Phase: Reconstruction}
\label{sec:online_exp}

With bases and responsibilities frozen, GatingNet and BranchNets are trained
by minimising~\eqref{eq:online_loss}.
Table~\ref{tab:results} reports test-set relative $L^2$ errors.
Figure~\ref{fig:reconstruct} shows a representative space--time prediction.

\begin{table}[h]
\centering
\caption{Relative $L^2$ test error on 1D Burgers' (mean over 1000 trajectories per~$\nu$).
  Each row of POD-DeepONet and NeuralPOD-DeepONet is a \emph{separate} model trained on
  the indicated~$\nu$; diagonal entries are in-distribution, off-diagonal show
  cross-$\nu$ generalization.
  TEMPO trains a \emph{single} model on all three $\nu$ simultaneously.
  \textbf{Bold} = best result in that test column.}
\label{tab:results}
\small
\begin{tabular}{llccc}
\toprule
Method & Trained on & $\nu{=}0.001$ & $\nu{=}0.1$ & $\nu{=}1.0$ \\
\midrule
\multirow{3}{*}{POD-DeepONet~\cite{Lu2022}}
  & $\nu{=}0.001$ & \textbf{0.2727} & 0.2970 & 0.6390 \\
  & $\nu{=}0.1$   & 0.3495 & \textbf{0.0480} & 0.3442 \\
  & $\nu{=}1.0$   & 0.4606 & 0.2597 & \textbf{0.0252} \\
\midrule
\multirow{3}{*}{NeuralPOD-DeepONet~\textit{(ours)}}
  & $\nu{=}0.001$ & 0.3915 & 0.2883 & 0.5351 \\
  & $\nu{=}0.1$   & 0.4259 & 0.2291 & 0.3714 \\
  & $\nu{=}1.0$   & 0.4828 & 0.3153 & 0.1830 \\
\midrule
\multirow{3}{*}{TEMPO(POD)~\textit{(ours)}}
  & $M{=}2$, all & 0.3187 & \textbf{0.1498} & \textbf{0.1590} \\
  & $M{=}3$, all & 0.3170 & 0.1677 & 0.1823 \\
  & $M{=}4$, all & 0.3215 & 0.1612 & 0.1998 \\
\bottomrule
\end{tabular}
\end{table}

\begin{figure}[h]
\centering
\includegraphics[width=\textwidth]{TEMPO/main/pictures/reconstruct.png}
\caption{Space--time reconstruction on 1D Burgers' ($\nu{=}1.0$):
  ground truth (\textit{left}), TEMPO(NeuralPOD) prediction (\textit{centre}),
  and point-wise absolute error (\textit{right}).}
\label{fig:reconstruct}
\end{figure}



\section{Conclusion}

Operator learning for parametric PDEs has largely assumed that a single learned representation suffices across the full parameter space.
When the governing parameter drives qualitative change in solution structure, this assumption forces the basis to compromise between incompatible geometries, concentrating error precisely where accurate prediction matters most.

TEMPO addresses this through a principled latent variable model.
Expectation-Maximization over POD coefficient space discovers $M$ latent regimes without parameter supervision, and each regime receives a dedicated locally adapted POD or NeuralPOD basis.
On 1D Burgers' equation spanning three orders of viscosity magnitude, the discovered regimes cross viscosity confirm that the partition reflects intrinsic solution geometry and  physical parameter label.
Regime-specific bases achieve substantially lower reconstruction error at equal mode count than their global counterparts, and the gated online operator outperforms POD-DeepONet and NeuralPOD-DeepONet across the full viscosity range, with the largest gains at extremes where solution character differs most.

Several directions remain open.
Conditioning the mode functions on the physical parameter would allow each basis to encode parameter-dependent solution structure directly, reducing the residual error that regime-agnostic modes cannot capture.
Extending TEMPO to multi-dimensional benchmarks will test whether regime discovery generalises beyond one-dimensional shock dynamics.
Convergence guarantees for Generalized EM under nonlinear basis parameterisation are a further theoretical open problem.

\newpage
\clearpage


\bibliographystyle{unsrtnat}
\bibliography{references}


\newpage
\appendix
\section{Appendix / supplemental material}\label{app}

\subsection{Additional experiments / Proofs of Theorems}\label{app:exp}



\end{document}